\documentclass[11pt]{article}
\usepackage[margin=1in]{geometry}
\usepackage{amsmath,amssymb,longtable,array}
\usepackage[hidelinks]{hyperref}
\setlength{\emergencystretch}{2em}
\newcommand{\mr}{\operatorname{mr}}
\newcommand{\F}{\mathbb F}
\title{Meridional rank conjecture literature and construction gates}
\author{Primary source research map}
\date{Evidence checked through 17 November 2026}
\begin{document}
\maketitle

The classical meridional rank conjecture remains open. It asks whether
$\mr(K)=b(K)$ for every knot $K\subset S^3$; the usual link formulation asks
the same equality for links. The 2026 problem list retains it as Problem
1.18 \cite{k3}. No general proof or classical counterexample was found in
this search. The catalogue \texttt{sources.json} records 48 sources,
their precise claimed scope and audit depth. This is a broad map of the
main results and recent developments, with outstanding source audits
retained rather than an assertion of exhaustive bibliographic coverage.

The graph manifold models now certify $\mr=b=4$ for all sixteen original
targets, whose complete list is in the final-five proof.
Together with the four earlier
seed cases, twenty of the retained293 names had equalities proved by
those methods. The arbitrary-cycle cover calculation below proves
another twenty-four, for forty-four retained named equalities. All sixteen original targets now have established
membership in the existing Boileau--Jang--Weidmann graph-cover MRC family,
using closed-manifold rigidity to identify their actual covers.
They are not cases outside known MRC families. Priority of the explicit
named bridge/rank calculations remains unestablished.
The sections below distinguish each method's scope from the current
sixteen-case conclusion. Full MRC remains open.

The curated checkout retains the current graph-cover proofs and their
listed finite inputs. Several earlier search panels and replay receipts
summarized in this literature map are available only through Git history;
those summaries do not provide a runnable replay of those older tests.

\section{What must be proved and how existing methods do it}

For any knot, a bridge presentation gives
\[
 \operatorname{rank}\pi_1(S^3\setminus K)\le\mr(K)\le b(K).
\]
The difficult direction is a meridian-sensitive lower bound as large as
a geometric bridge upper bound. Ordinary rank is insufficient: nontrivial
torus knot groups have rank two while their bridge numbers are unbounded.
Normal generation by one meridian is also different from generation by
one meridian.

If $G_K\twoheadrightarrow Q$ maps meridians into a distinguished conjugacy
class, every meridional generating set induces a generating set in that
class. Coxeter reflections, symmetric-group transpositions and affine
transvections exploit this restriction. Knot contact homology instead
supplies augmentation-rank bounds. A finite quotient need not preserve the
preferred longitude unless the particular application requires it; the
meridian conjugacy data are essential in all of these lower bounds.

The geometric upper bound can be exhibited by a bridge diagram or by a
diagram with a specified Wirtinger seed set and a valid propagation
sequence. Blair--Kjuchukova--Velazquez--Villanueva prove that minimum
diagram Wirtinger number equals bridge number \cite{wirt}. This does not
identify arbitrary meridional generators with diagram seeds.

\section{Established families and their actual scope}

\begin{enumerate}
\item \textbf{Low ranks.} Boileau--Zimmermann show that meridional rank two
characterizes two-bridge links. Rank one characterizes the unknot in the
knot case. Thus all knots of bridge number at most three satisfy MRC:
if a three-bridge knot had meridional rank at most two, it would have
bridge number at most two. The reverse implication
$\mr(K)=3\Rightarrow b(K)=3$ for an arbitrary knot is a different,
unsettled assertion. The low-rank statements are restated in the primary
introduction of \cite{bjw}.

\item \textbf{Torus and Montesinos results.} Rost--Zieschang prove
$\mr(T(p,q))=b(T(p,q))=\min(|p|,|q|)$ for torus links \cite{rz}.
Boileau--Zieschang prove the Montesinos results. Lustig--Moriah's
generalized Montesinos theorem applies when the gcd $\alpha$ of the
constituent two-bridge numerators differs from one: for their $2n$-plat,
ordinary group rank, exterior Heegaard genus and bridge number all
equal $n$, hence so does meridional rank \cite{moriah}.
This numerical scope is extracted from the coauthor's survey; the
original paper's full proof and geometric definition remain unaudited.
Pretzels with at least three parameters all of absolute value
at least two are explicitly recovered by the cycle-quotient construction
in \cite{bbk}. Burde's 1988 dihedral-cover paper is retained as a historical
source in the catalogue; its introduction alone is not used as a broader
MRC theorem.

\item \textbf{Graph manifolds and controlled satellite operations.}
Boileau--Dutra--Jang--Weidmann prove MRC for the class built from
meridionally tame knots by connected sums and braid-pattern satellites;
all knots with graph-manifold exteriors are included \cite{bdjw}.
This does not make arbitrary satellite patterns or arbitrary connected
sums of arbitrary MRC knots automatically admissible.

\item \textbf{Double branched covers and prime satellites.}
Boileau--Jang--Weidmann show that a link with meridional rank three and
graph-manifold double branched cover has bridge number three. A
four-bridge link in that cover class has meridional rank four.
For a prime knot with nontrivial JSJ decomposition they prove
$\mr\ge4$ \cite{bjw}. The prime and nontrivial-JSJ hypotheses avoid an
incorrect blanket claim for composites.

\item \textbf{Augmentations.} Cornwell develops representation and
augmentation lower bounds, with applications to torus knots and a
three-bridge pretzel family \cite{cornwell}. Cornwell--Hemminger prove a
multiplication theorem for specified braid-pattern satellites and obtain
large iterated torus-knot classes \cite{ch}. The pattern and maximal-rank
hypotheses are material; this is not a theorem for all satellites.

\item \textbf{Coxeter, twisted and arborescent links.}
Baader--Blair--Kjuchukova prove MRC for twisted links and for links from
their even-weight bipartite trees \cite{bbk}. Their use of ``bipartite''
requires all branching vertices to have the same colour, not merely
the ordinary fact that every tree is bipartite.
Baader--Blair--Kjuchukova--Misev prove $\mr=b=m(T)$ for a weighted plane
tree $T$ obtained from a subtree $T_0$ by adding at least three twigs
at every vertex of $T_0$, provided every weight is outside
$\{0,1,-1\}$. A twig is a straight branch from a leaf to a branch
point. Their $m(T)$ equals $f(T)+2$, where $f(T)$ is the minimum
number of edges whose deletion leaves a forest of maximum valence two
\cite[Theorems 1 and 2]{twigs}. The three-twig condition is material.
Neither result states MRC for every arborescent link.
Baader--Kjuchukova's symmetric quotients also prove MRC for an
unknotting-number-one family of unbounded bridge number \cite{bk}.

\item \textbf{Finite computations.} Blair--Kjuchukova--Morrison find
explicit maximal-rank Coxeter quotients for 595,515 of 1,701,936 knots
through 16 crossings \cite{bkm}. This is not a complete verification
through 16 crossings, nor a count of everything known by other methods.
Their finite quotient limitations are relevant to the proposed searches.
\end{enumerate}

Lustig--Moriah's broader \emph{plat} criterion can be applied without
assuming an unspecified family label. For a $2n$-plat braid $B$, let
$\widehat\rho(B)$ be its unreduced Burau matrix evaluated at $t=-1$.
Let $L$ have row $i$ equal to $e_{2i-1}^{\mathsf T}-e_{2i}^{\mathsf T}$,
and let column $i$ of $M$ be $e_{2i-1}+e_{2i}$. If the gcd of all
entries of $L\widehat\rho(B)M$ differs from one, then ordinary group
rank, exterior genus and bridge number equal $n$ \cite[Theorem 4.20]{moriah}.
The extra condition on a determinant $\beta$ concerns inequivalent
tunnel systems, and is not required for this equality. Applying this
test to a new candidate still requires its explicit plat and the
source's conventions; no claim is made for the gcd-one case.

\section{Recent work that changes the research map}

\textbf{Fishnets and stabilization, published 2026.}
Blair--Kjuchukova--Pfaff give fishnet families satisfying MRC with
unbounded bridge number and arbitrarily high distance relative to bridge
number. They also show that, for a $b$-bridge position with $b>1$, a
suitable unknot $U$ can be added so that
$b(L\cup U)=\mr(L\cup U)=\operatorname{rank}G_{L\cup U}=2b-1$
\cite{fishnet}. The added link's equality does not imply equality for the
original sublink. Their separate plain-sphere presentation invariant equals
bridge number, but retains a restriction on permitted relations \cite{plain}.

\textbf{Whitehead doubles, November 2026 preprint.}
Gao--Gu--Li--Wang claim that for every nontrivial two-bridge knot and every
$N\ge0$, its $N$-fold untwisted Whitehead double has
$\operatorname{ar}=\mr=b=2^{N+1}$. They use microlocal sheaves and iterable
peripheral curves. They also give a square-knot double with augmentation
rank at most five but $\mr=b=6$. Their introduction identifies a specific
gap in Dutra's earlier Whitehead-double preprint; it does not disprove
its conclusion. The new 44-page proof has not been independently audited
here \cite{whitehead}. Thus Dutra's broader claimed theorem remains
separately flagged, rather than silently counted as established.
The original v1 pages seven and twelve are now directly checked:
$\alpha(y_1y_2)=\omega(y_1y_2)=x_2x_1^{-1}x_2^{-1}x_1$ is a
nontrivial freely reduced word in both subgroups, contradicting
Lemma 3.1 at conjugator one. Lemma 3.8(2) explicitly invokes it.
No proved record in this programme depends on the 2022 theorem;
the earlier published graph-manifold result is a distinct source.

\textbf{Transvection computations, October 2026 deposit.}
Bhattacharjee's unrefereed release claims that 243 of 536 bridge-uncertain
13-crossing prime knots have $b=\mr=4$, using symplectic transvections in
characteristic three or five, or orthogonal reflections. The deposit
leaves 293 in that study and includes verifiers \cite{deposit}. All 243
algebraic and four-seed certificates were reported to pass a separate
standard-library replay. Its receipt is not distributed in this curated
checkout. The same audit reported that all 243 named diagrams match the
independently packaged HT census
by exact planar-map bijections. The author's PD data agree literally
with KnotInfo snapshots 2026.10.5 and 2026.11.2. Both snapshots still list
the same 536 cases as $[3,4]$; this table state does not negate the
certificate proofs. The full paper and publication priority remain
unaudited; the release is unrefereed.
The odd-characteristic linear method
does not satisfy our binary affine transitivity hypotheses automatically.

\section{What the affine argument now supplies}

The companion \texttt{PROOF.tex} supplies three direct bridges:
\begin{enumerate}
\item The original $k+1$ lower bound extends to every \emph{nonzero}
common alternating form, including degenerate forms. This permits both
parities of bridge number.
\item Explicit representations of
$P(2m,3(2h_1+1),\ldots,3(2h_{r-1}+1))$ onto
$W(D_r)=\F_2^{r-1}\rtimes S_r$ act transitively on $2^{r-1}$ sheets.
The standard bridge diagram proves the matching upper bound $r$.
For odd $r$ these meet the original affine-symplectic hypotheses exactly.
These are already known MRC cases, so they establish applicability rather
than new knot-table results.
\item Given a specified transposition quotient onto $S_r$, a finite
binary linear system decides whether it lifts to a transitive signed
action. If its solution-space dimension exceeds $r-1$, the lift proves
$\mr\ge r$, one more than the underlying symmetric quotient bound.
\end{enumerate}

The last test is especially useful when a known $S_r$ colouring gives only
$r-1$ while a bridge diagram gives $r$. Both the presentation's knot
identity and meridian labels, and the geometric upper certificate, must
be established independently. No such new unresolved example is claimed
in the present work.

The companion \texttt{ALL\_CHARACTERISTIC\_LIFTS.tex} extends the intrinsic
lift-dimension bound to every field. A single integer relator matrix and
its certified Smith form give the maximum dimension over all
characteristics. For each of sixteen frozen five-seed table diagrams,
an $S_4$ colouring has lift dimension five over $\F_5$, proving
$\mr=b=5$. Separate named-census transfer checks all sixteen certificates
directly on their diagrams. An independent rational-tangle constructor
and exact planar-map bijections identify every case as a five-tangle
Montesinos knot. Thus all sixteen equalities were already covered by
Boileau--Zieschang (1985), whose Theorem 1.1 is explicitly restated in
Musick's Theorem 1 \cite{musick}. Half-turn diagram conventions are
justified by the rational-tangle flipping lemma \cite{kl}. The earlier
binary bounds remain four. These are useful controls, not new MRC cases.

This counting method has relevant quandle precedents:
Clark--Elhamdadi--Saito--Yeatman define fibre counts of colouring lifts
and use total colouring counts for bridge bounds \cite{cesy}. Our fixed
group-representation lift space is such a fibre, and evaluation on
arbitrary meridional generators supplies its meridional lower bound.
This connection is now part of the priority audit.

The companion \texttt{BRANCHED\_COVER\_LIFT\_BOUNDS.tex} identifies the
integer lift matrix with the first homology of the simple degree-$r$
branched cover $Y_\phi$ associated to a transposition quotient:
\[
 \dim_F Z_\phi(F)=r-1+\dim_F H_1(Y_\phi;F),\qquad
 \mr(K)\ge r-1+\operatorname{rank}\pi_1(Y_\phi).
\]
Its nonunit Smith factors are exactly the torsion invariant factors of
$H_1(Y_\phi;\mathbb Z)$. A finite nonabelian coefficient group $A$
gives $|A|^{r-1}|\operatorname{Hom}(\pi_1Y_\phi,A)|$ admissible wreath
lifts. For $r=2$ this recovers the classical core group
$\pi_1(\Sigma_2K)*\mathbb Z$, due to Kelly and Wada, as reviewed and
developed by Silver--Traldi--Williams \cite{core}. No priority claim is
made for the cover-rank or counting framework. The uniform arguments
are working proofs without formal verification or external review.

On the frozen sixteen current thirteen-crossing targets, four-seed
upper certificates force $r\le5$ for a transposition quotient.
Complete searches in degrees two through five, followed by integer
Smith replay, show that the symmetric root-lift method gives at most
two in \emph{every} degree and characteristic. The original binary
criterion is impossible there in every dimension: $k+1\le b\le4$
forces $k\le3$, hence a rank-two form and the absent $S_3$ quotient.
The previous suggestion of higher-rank binary forms on this panel
is corrected. A separate $A_5$ core-colouring test, reproduced with
$\operatorname{PSL}_2(\F_5)$ matrices, also gives no sharp four-bound.
Other quotient and nonabelian cover methods remain open.

The user's joint-quotient proposal adds a stronger cover-rank test:
twenty distinct $A_5$ quotient kernels of one double-cover group
force rank at least three and $\mr\ge4$. Burness--Crestani explicitly
give the classical capacity nineteen, crediting Hall's direct-power
formula \cite[p.~2]{bc}. The quotient kernels must have a single
domain; witnesses on different knots cannot be pooled. The companion
\path{JOINT_SIMPLE_QUOTIENT_KERNELS.tex} proves the capacity bound
and tests all 293 retained uncertainty names, with maximum five and
no sharp case. The known control $P(2,3,3,5)$ has 23 kernels and
2901 homomorphisms, so the joint test gives sharp four while the raw
count gives only three. The Seifert presentation in
\cite[Section 7.3]{ps} also gives a uniform known family with
denominators divisible by $2,3,3,5$. These remain established
Montesinos cases. The original Hall paper and full Burness--Crestani
and Poudel--Saveliev articles have not been independently audited;
the cited formula and Seifert presentation were read at their stated
locations and retained with hashes.

The user's augmentation-module refinement uses an index-$n$ subgroup
$H$ of a branched-cover group $G$, in characteristic prime to $n$.
Shapiro gives $\Delta=b_1(H)-b_1(G)=\dim H^1(G,V_{n-1})$, and
$\operatorname{rank}G\ge1+\lceil\Delta/(n-1)\rceil$. For degree five in the
double cover, gain five implies $\mr\ge4$. All sixteen $A_5$
kernel classes on the retained first-sixteen panel give gain at most
one in every prime. Their base abelianizations, coprime to thirty,
force every transitive degree-five action to have image $A_5$,
so all connected index-five covers are covered by this negative
test. The proof, direct cocycle replay and sharp known controls are
in \path{AUGMENTATION_COVER_HOMOLOGY.tex}. This uses a different
group and module from the old fixed knot-root exclusions.

The same note records a concrete gate for the peripheral direction:
in a representation preserving a nondegenerate symmetric form, a
diagonalizable rank-one meridian has an anisotropic eigenline.
The commuting longitude acts on it with $\lambda^2=1$, the excluded
peripheral locus in \cite[Section 2.3]{whitehead}. A useful deformation
must escape that constraint and supply actual augmentation data;
arbitrary peripheral parameters do not suffice. No such curve on
the sixteen targets is constructed here.

\section{Obstructions and questions that remain genuinely open}

The 2026 problem list retains finite-quotient detection of meridional
rank, bridge monotonicity under meridian-preserving epimorphisms, and
freeness of subgroups on fewer than $b(K)$ meridians for hyperbolic knots
as related questions \cite{k3}. A meridian-preserving epimorphism
$G_{K_1}\twoheadrightarrow G_{K_2}$ does give
$\mr(K_1)\ge\mr(K_2)$ immediately, but bridge monotonicity would require
more. General connected-sum additivity
$\mr(K\mathbin{\#}J)=\mr(K)+\mr(J)-1$ is also not settled by proving it
for quotient-detected families.

Every meridian-involution representation factors through
$G_K/\langle\!\langle\mu^2\rangle\!\rangle$, the $\pi$-orbifold group.
Boileau--Kitano--Nozaki explicitly formulate the stronger equality between
bridge number and torsion-generator rank of the orbifold group for prime
links with infinite orbifold group \cite{preorder}. Lifting a small
generating set of the quotient need not generate the knot group.

Finite Coxeter quotients do not universally detect meridional rank.
For $12a210$, an infinite triangle Coxeter quotient is sharp while no
finite Coxeter quotient is sharp. Ryffel's theorem, reproduced in
\cite{bkm}, gives odd torus knots with no nontrivial Coxeter quotient at
all. Consequently a universal search restricted to $W(D_r)$ cannot solve
MRC. The companion proof also gives a stronger elementary limitation:
for $p\ge3$ and prime $q>2^{p-1}-1$, $T(p,q)$ has no representation
satisfying our common-form affine criterion in any dimension. The bound
would force $k\le p-1$, while the image order in $\operatorname{AGL}(k,2)$
is prime to $q$. The central torus relation then makes the image abelian,
of order at most two. Thus even allowing non-Coxeter affine images
does not make this criterion universal.

The companion quadratic-coefficient note extends exact field testing
to $\mathbb Z[\tau]$, $\tau^2=\tau+1$, using expanded integer Smith
matrices and separate residue embeddings at rank-drop primes. It
uses twelve existing first $A_5$ core maps to obtain $H_3$ reflection
bases. Every fixed base has lift dimension three in every field.
An integral split pullback shows that all these affine lifts factor
through the same 120-element group. In characteristic two, a retained
lift has no common fixed point yet has four orbits of sizes
$12,12,20,20$ on $\mathbb F_4^3$. Thus the fixed-base affine route is
closed without asserting global fixed points or excluding other
base colourings. A direct comparison with \cite[Lemma 6.2]{whitehead},
whose complete elementary proof has now been read, shows that the
unrestricted cocycle estimate cannot improve the restricted root-lift
bound on the same representation.
The original nondegenerate affine-symplectic $k+1$ bound also follows
by combining its $k$-dimensional principal cocycles with the
nonprincipal cocycle of the transitive action. The degenerate-form
extension still requires the retained quadratic proof. Details are in
\path{ICOSAHEDRAL_LIFTS_AND_QUADRATIC_COEFFICIENTS.tex}.

Further work needs a structurally motivated cover degree or an actual
augmentation curve outside the orthogonal locus. The outstanding
parameter and recent-proof audits are listed in \texttt{sources.json}.

\section{Formal deformation gate at icosahedral points}

The companion working proof
\path{SIX_POINT_COVERS_AND_FORMAL_ORTHOGONALITY.tex} tests the
different six-point $A_5$ action on all sixteen retained kernel
classes. Subgroup homology is cyclic of order $4D$ or $28D$,
giving no sharp bound. The six-point permutation character is the
symmetric square of either standard complex three-dimensional
icosahedral character. For $G=\pi_1(\Sigma_2K)$ and its index-six
subgroup $H_6$, Shapiro and index-two restriction from the knot
orbifold group $Q$ therefore give
\[
 H^1(G,\operatorname{Sym}^2W^*)\cong H^1(H_6;\mathbb C)=0,
 \qquad H^1(Q,\operatorname{Sym}^2W^*)=0.
\]
The form-lifting proof uses this vanishing at the original point:
at each formal order, failure of form invariance is a cocycle and
can be removed by a correction. Thus every formal complex
deformation through these points that keeps meridians involutive
preserves a possibly varying nondegenerate symmetric form. Its
longitude still has $\lambda^2=1$ on the distinguished meridian
eigenline. A family with varying meridian eigenvalue need not factor
through $Q$ and remains outside this obstruction. These are internal
working proofs, unformalized and externally unreviewed.

Separate direct cocycle checks at the nonsplit small primes complete
the earlier field boundary: the ordinary cover-generation cocycle
estimates from both permutation augmentation modules fail in every
characteristic. A frozen primitive-word probe found no two-generator
certificate. Neither these computations nor the formal-form argument
prove an actual cover-rank upper bound or settle a new MRC case.

\clearpage
\section{Actual peripheral curves for torus Whitehead iterates}

The companion working proof
\path{TORUS_WHITEHEAD_PERIPHERAL_CURVES.tex} supplies a positive
application of varying meridian eigenvalues. For coprime
$2\le p<q$, an explicit companion-matrix family of irreducible
rank-$p$ KCH representations of $T(p,q)$ realizes
\[
 \lambda=(-1)^{p+1}\mu^{-(p-1)q},\qquad
 \eta=\lambda^2=\mu^{-2(p-1)q}.
\]
The maximal torus branch is already in Cornwell's work
\cite[Section~5.1]{cornwell}. Here generic irreducibility is proved
directly by a cyclic-vector argument for the two companion matrices,
so the branch consists of actual rank-$p$ points away from a finite
set. At $\mu=\infty$ its peripheral orders are $(-1,2(p-1)q)$.
They meet the iterable-branch hypothesis of
\cite[Definition~4.1 and Corollary~4.5]{whitehead}, giving the
source-derived corollary
\[
 \operatorname{ar}(\operatorname{Wh}^N(T(p,q)))
 =\mr(\operatorname{Wh}^N(T(p,q)))
 =b(\operatorname{Wh}^N(T(p,q)))=p\,2^N.
\]
The fixed untwisted pattern and conventions are those of that
preprint; this does not use the affected Dutra 2022 lemma.

There is a broader sufficient criterion: if an $n$-strand braid
closure has a rank-$n$ augmentation and nonzero writhe $w$, all
its iterates have the three ranks $n2^N$. Cornwell's
Theorem~1.10 \cite{cornwell17}, including the last paragraph of its
proof, gives actual rank-$n$ points for all but finitely many
$\mu$ on $\lambda=(-\mu)^{-w}$. Its squared longitude has orders
$(-1,2w)$, satisfying the same gate for either sign of nonzero $w$.

For $T(3,5)$, the explicit family has $\mu=s^3$ and
$\lambda=s^{-30}$. At $s=-1$ its image is the 120-element
icosahedral reflection group. Both eigenvalues vary there, giving
an actual escape from the fixed-involution formal-form barrier.
This is a torus control, not a representation curve for the sixteen
table targets.

The required one-step algebra and iterable valuation argument have
now been read and checked; exact symbolic block identities and
frozen finite controls pass. The geometric peripheral inputs and
the preprint theorem remain explicit dependencies. Its full
microlocal theory, two-bridge existence proof and other sections
are not independently certified. This is a derived application
without a priority claim, formal verification, full MRC closure
or a newly settled member of the retained panel.

\clearpage
\section{Varying meridians and the stabilization boundary}

The companion note
\path{MERIDIAN_DEFORMATIONS_AND_STABILIZATION.tex} now tests every
retained icosahedral base with the meridian eigenvalue allowed to
vary. The sixteen bases occur on twelve target knots; the four
targets without an $A_5$ quotient supply no such point.
Each normalized tangent matrix has size $144$ by $28$, rank $23$
and kernel dimension five. Four dimensions are residual
conjugations and the remaining direction has nonzero meridian
eigenvalue differential. Exact quadratic-rational kernel bases and
rank lower certificates have separate dual-number/minor replay.

The fixed-eigenvalue kernel consists only of conjugations. An
order-by-order conjugation argument therefore strengthens the
previous invariant-form result: every formal rank-three deformation
at these points that keeps meridians involutive is conjugate to the
constant representation. Allowing the eigenvalue to vary leaves
one tangent direction. A transverse slice and the analytic implicit
function theorem reduce its integrability to residual functions of
one parameter. They must vanish identically to give an actual curve;
a finite nonzero obstruction instead gives an isolated thick point.

There is also a uniform result for adjoining trivial summands. If
the restricted first cohomology vanishes for both $W$ and $W^*$,
every formal rank-one-meridian family through $\rho_0\oplus I_d$
splits as $\sigma_t\oplus I_d$. Its meridian eigenvalue may vary.
At the first nonzero off-diagonal order, rank-one minors force
admissible cocycles; their vanishing modulo principal cocycles
removes both blocks by conjugation. The trivial block is forced to
remain trivial. Both restricted cohomology spaces vanish on all
sixteen retained bases, by independent rational/affine replay.
Thus adjoining any number of trivial summands at these points
cannot produce a sharp higher-rank irreducible representation.

All $32$ frozen degree-three jets, using both split embeddings
modulo $11$, satisfy every oriented relator and all rank-one
minors. These residue-field jets do not certify a complex curve or
all-order integrability. The working proofs are unformalized and
externally unreviewed, with no novelty or new MRC case claimed.
The remaining choices include proving an actual rank-three curve,
finding a higher-rank family through another special fibre, or
obtaining another certified cover-rank bound. Full MRC stays open.

\clearpage
\section{Complete ordinary coefficient boundary for a single A5 map}

The working proof
\path{SIMPLE_MODULE_BOUNDS_AND_A5_EXHAUSTION.tex} completes the
ordinary cover-group cocycle test through the retained $A_5$ maps.
An elementary reduction makes this finite: in an exact sequence
$0\to U\to V\to W\to0$,
\[
 \dim Z^1(G,V)\le\dim Z^1(G,U)+\dim Z^1(G,W).
\]
The largest coefficient-normalized dimension is therefore attained
on a simple module. This includes nonsemisimple extensions and
commutes with scalar extension. The absolute simple-module lists
are imported from the primary ATLAS representation page \cite{atlas}.

The previously untested three-dimensional modules now have exact
integral matrices on all sixteen quotient-kernel classes. Their
regular expansions have nonunit Smith factors $5$ on eight classes,
and $7,35$ on eight others. Separate treatment of the ramified prime
five and every residue embedding completes the all-field check.
The two-dimensional characteristic-two quotient modules and their
Frobenius conjugates have cocycle dimension three. Modular
matrix-algebra spans identify all supplied simple models. Combined
with the earlier complete $V_4,V_5$ and trivial-module data, this gives
\[
 \dim Z^1(G,V)\le\tfrac32\dim V
\]
for every finite-dimensional coefficient module through any single
$A_5$ map on the sixteen-target panel, over every field. Thus none
of those unrestricted linear estimates can force ordinary cover
rank three. This does not prove that the groups are two-generated.

Shapiro also gives an exhaustive consequence for induced covers:
for $H=\phi^{-1}(L)$, $L\le A_5$, and $n=[G:H]$,
\[
 b_1(H;k)\le\lfloor n/2\rfloor+1.
\]
All such covers have zero rational first Betti number, including
the degree-sixty kernel covers. Their Schreier and split
augmentation-module gain tests cannot reach the rank-three gate.
This concerns all $A_5$-induced degrees, not other monodromy.

Two retained kernels on the same knot give a different joint
quotient. Its natural four-dimensional tensor modules over
$\F_4$ have been checked on the four eligible knots: all sixteen
Frobenius choices have $\dim Z^1=6$ and $\dim H^1=2$, below the
threshold nine and five, respectively. Other product modules and
quotients remain open. Independent integral/affine, Smith, quotient
and matrix-algebra replays pass. No new knot equality or full MRC
claim follows; uniform proofs remain unformalized and unreviewed.

\section{Singular affine certificates and the full GL3 quotient panel}

The singular affine construction below gives a positive known-family
certificate. The reported negative finite-panel tests are separate from
that written construction; their former companion proof and complete
search inputs are not distributed in this curated checkout. Homogenizing a three-dimensional affine representation
can give rank-one meridians and zero common fixed space in dimension
four, even though every image direction lies in a three-space. A
nonprincipal meridian-admissible cocycle supplies precisely this
situation when the linear action has no fixed vectors.

Explicit binary matrices give sharp certificates for every
$T(4,7(2r+1))$, $r\ge0$. The torus presentation and meridian formula
are standard \cite{cornwell}; the MRC equalities were already known
\cite{rz}. The $T(4,7)$ control has image $AGL_3(\F_2)$ of order
1344, transitive on eight affine points. Its linear image preserves
no nonzero alternating form. For $3\nmid7(2r+1)$ the absence of
Fox three-colouring, together with the retained rank-two theorem,
excludes every realization of the original common-form criterion.
This comparison has no novelty claim for the knot equalities.

The source audit of the external binary search reports that it assumes four independent seed image directions and
discards singular pairing matrices. All homogeneous affine
extensions just described have pairing rank at most three, for
every four-meridian choice. They can satisfy the valid rank-one
fixed-space lower-bound criterion while lying outside that search.
This is a limitation of search coverage, not a defect in any of
the 243 deposited certificates. The separate three-seed program
uses bounded random diagram moves, not an exhaustive geometric
three-bridge test. The archived open-list bridge intervals are
$[3,4]$, so bridge number four is not an input assumption.

Motivated by this missing architecture, all 293 pinned open names
were frozen and exhaustively tested for double-cover quotient
kernels with target $GL_3(\F_2)=PSL_2(\F_7)$. Independent projective
replay gives 57 generating-pair and 13,368 generating-triple
automorphism orbits. The kernel-count histogram is
$139,123,29,1,1$ knots with respectively $0,1,2,3,7$ kernels.
The same-source rank-three threshold would be 58, so there is no
sharp kernel-count case. All 191 maps on 154 knots have binary
natural/dual cocycle dimension four or five and Steinberg dimension
eight or ten. The characteristic-two absolute simple list
\cite{atlasgl}, exact matrix-algebra checks and composition-series
argument give $\dim Z^1(G,V)\le5\dim V/3$ for every module through
each epimorphism. This includes nonsplit extensions over arbitrary
characteristic-two fields. It does not cover all proper-image
representations on the enlarged panel or other characteristics.

The 29 inner deck maps give 58 natural/dual binary meridian bases,
all with root-lift dimension three. The sign-twisted adjoint
construction over $\F_7$ covers all 191 maps of either deck type;
all root-lift spaces again have dimension three. The independent
verifier checks group identifications, full tuple orbits, every
diagram closure, coefficient matrices and nullspaces, as well as
the positive torus control. Source PD parsing is shared; uniform
proofs remain unformalized and externally unreviewed. No new open
knot is closed. Full MRC, actual target group ranks and broader
source/priority audits remain open.

\section{Exact scope of the rank-two-form method}

The new working proof \path{RANK_TWO_AFFINE_EQUIVALENCE.tex}
gives a converse construction for the original method, including
arbitrarily large radicals, when the common form has rank two.
Fix an onto $S_3$ transposition colouring and its simple
three-sheeted branched cover $Y$. A transitive certificate in
dimension $k\ge2$, projecting to that colouring, exists exactly
when $b_1(Y;\F_2)\ge k-1$. Its maximal dimension is $b_1+1$ and
its best meridional bound is $b_1+2$, when $b_1>0$. If $b_1=0$
no such representation exists. This identifies the exact gain
and a computable stopping condition for this part of the
unpublished quadratic criterion; it is not a universal detection
theorem or a novelty claim.

The proof first establishes a general knot-group lemma:
quotienting a meridian-rank-one module by a trivial submodule
preserves its admissible cocycle space when the image lines
survive. Vanishing $H^2(\Gamma;R)$ and evaluation
$H^1(\Gamma;R)=R$ at a meridian lift and normalize the cocycle.
Both facts follow from the knot exterior's homology and the
universal coefficient theorem. Dependent radical/translation
classes preserve an explicit nonconstant quadratic function.
Independent classes give the full affine image of order
$24\cdot8^{k-2}$, with transitivity proved by restriction to
the $S_3$ kernel and central commutators.

There are exactly $2^{b_1(Y;\F_2)}-1$ distinct $S_4$
transposition quotient kernels above this one $S_3$ kernel.
Thus two distinct such kernels on the same knot group already
force meridional rank at least four and construct an eight-point
rank-two-form certificate. The common lower quotient is essential;
no pooling of knots or unrelated quotients is allowed. For an
$S_4$ colouring, the induced covers also satisfy
$b_1(Y_3;\F_2)=1+b_1(Y_4;\F_2)$ by the radical-lifting lemma.

The complete 293-name symmetric scalar audit gives 26 $S_3$
bases, with binary lift dimension two on five and three on 21;
there are 21 $S_4$ bases and no $S_5$ base. All 170 retained
$A_5$ cover kernels have inner deck action, independently
agreeing with the absence of an outer extension to $S_5$.
Every degree and characteristic has scalar bound at most three.
Since the original quadratic bound itself forces $k\le3$ on
these upper-four targets, rank two is the only possible nonzero
form rank. Exactly 21 targets have a dimension-two realization,
giving bound three; the other 272 have no original representation
in any dimension. No target has a sharp original certificate.

The proper-image gap in binary three-dimensional affine bases
is now closed as well. Complete subgroup and normal-generation
replay leaves $C_2,S_3,S_4,GL_3(\F_2)$; the first three reduce
to the new tests, while the retained 58 full-image binary root
tests are reused. Every such base has admissible dimension at
most three. This excludes the specified homogeneous affine
architecture on all 293 knots, not other coefficient fields,
larger dimensions or unrelated representations.

Three actual known pretzel controls realize dimensions two,
three and four, with affine image orders 24, 192 and 1536.
Their independent-class constructions are transitive; dependent
classes preserve the predicted quadratic invariant. Complete
tuple orbit, diagram, affine-matrix, Smith and finite group
replays pass, with source PD parsing and older replay routines
shared. Uniform proofs remain unformalized and externally
unreviewed. These controls do not close a new open knot; actual
target ranks, full MRC and broader source/priority audits remain
open.

The initial nonlinear follow-up left seven explicit
two-generator cover presentations after proving 19 cyclic.
The small-quotient and bounded-coset records remain as that
checkpoint. The structural continuation below now resolves
all seven and produces four sharp knot certificates; it does
not change the preceding scalar or original-criterion
exclusions.

\section{Four sharp cases from nonabelian three-sheeted cover groups}

The working proof \path{TREFOIL_AMALGAMS_SHARP_FOUR.tex}
establishes
\[
 \mr(K)=b(K)=4\quad\text{for}\quad
 K=13a_{3145},13a_{3147},13a_{3148},13a_{3149}.
\]
These names lie outside the original sixteen-target prefix
and remain in the external project's 293-name open list
fetched on 17 November. They are absent from its deposited
243 certificates. Worldwide novelty and exclusion from every
known family remain unaudited, so a previously-unresolved
or priority claim is withheld.

Exact presentation isomorphisms, including a targeted
invertible basis change for $13a_{3148}$, give
\begin{align*}
 A_\epsilon&=\langle a,b\mid b^3a^2b^3=a,
                  \ a^{3\epsilon}ba^{3\epsilon}=b^2\rangle,\\
 B_\epsilon&=\langle a,b\mid b^2ab^2=a^2,
                  \ (a^{3\epsilon}b)^3=b^2\rangle.
\end{align*}
The four covers have groups $A_+,A_-,B_+,B_-$ respectively.
Each is an amalgam of two trefoil groups along the
meridian/fibre $\mathbb Z^2$ subgroups. The fibre has infinite order;
modulo it, the meridian has infinite order in $C_2*C_3$,
proving peripheral injectivity. The classical amalgam
normal-form theorem \cite{serre} embeds the nonabelian vertex
group. Thus every cover group has rank exactly two.
The general meridional cover bound gives $\mr\ge2+2=4$,
and independently replayed four-seed diagrams give $b\le4$
through the published Wirtinger-number theorem \cite{wirt}.

The named PDs match literally in both archived KnotInfo
releases and an independent HT census, after the standard
edge-label shift. A separate planar under-to-under arc
trace rebuilds all branch words and upper certificates;
the producer's parser is used only to transport indices.
Every relator shortening, inverse Nielsen map and Tietze
elimination is replayed. The conceptual peripheral and
normal-form arguments are written proofs, unformalized and
externally unreviewed, rather than inferred from finite tests.

There is now also a finite nonabelianity proof for these same
four covers: explicit pairs in $S_{25}$ satisfy all six literal
two-generator checkpoint relators and have nontrivial commutator.
All 39 original branched-cover relators per knot also pass an
independent pointwise replay. On $\mathbb F_5^2$, a trefoil action
has generators $\alpha(x,y)=(p(x),y+2)$ and
$\delta(x,y)=(q(x),y+3)$, with image lists
$p=(3,0,2,1,4)$ and $q=(4,1,3,2,0)$.
The fibre and meridian become independent translations, so
linear changes of coordinates realize the required peripheral
gluings. This produces the four 25-point witnesses directly.
Their vertex image is the nonsolvable group $A_5\times C_5$,
consistent with the solvable-quotient exclusion below.
The literal words, permutation arrays and standalone
standard-library checker are retained as
\path{finite_cover_certificates.json} and
\path{check_finite_cover_certificates.py}.
The full replay rebuilds the branch words from independent
planar arc traces and rejects four corrupted certificates.
This finite proof does not require amalgam normal form or
residual finiteness, and supplies no additional named case
beyond the four already retained.

The other three unresolved checkpoint groups are now
proved cyclic of orders two, two and four. All 26 tested
three-sheet cover-group ranks are therefore determined:
nineteen groups $C_2$, two $C_4$, one trivial, and the four
rank-two amalgams. The four latter groups have cyclic
abelianizations of orders 37,35,37,35, and their universal
abelian covers have zero integral first homology, by exact
Schreier matrices and replayed unit Smith forms. Their
commutator subgroups are nontrivial and perfect, so every
solvable quotient is cyclic. This explains why enlarging
the earlier solvable searches could not have found the
rank gain. All symmetric scalar lift bounds still fail
sharp four, and the original quadratic criterion has no
representation on these four knots. The nonabelian group
structure supplies the missing information.

The result gives four explicit named MRC certificates in
this programme. It does not settle the original sixteen
targets, the other 289 names of the frozen project list,
full classical MRC, or the outstanding literature/priority
audit. One stale duplicate many-twigs catalogue entry has
been reconciled with its fully scoped retained entry;
Serre's standard amalgam background is now separately
credited. That checkpoint had 35 distinct source identifiers.
Those two scoped primary-source additions brought the catalogue to 37;
the connected-sum sources below bring the current total to 39.

\section{Finite block actions and periodic sharp-four families}

The working proof \path{PERIODIC_TWIST_FAMILIES.tex} extends
each of the four retained seed cases to an infinite family.
Induce its 25-point cover action to a 75-point knot action
on three blocks. Meridian images exchange two blocks and
fix the third pointwise. If $d$ such exchanges generate a
transitive action on $r$ blocks, identifying fibres along
a spanning tree makes all tree exchanges trivial on the
labels. Only $d-r+1$ non-tree edge permutations remain to
generate a block-stabilizer action. A nonabelian stabilizer
therefore forces $d\ge r+1$, giving meridional rank at least
four in the retained three-block actions.

At the first crossing of every seed, the incoming meridian
images $x,y$ satisfy $(xy)^3=1$. The Hurwitz move
$T(x,y)=(y,yxy)$ has $T^2$ equal to conjugation by $(xy)^{-1}$,
so $T^6=1$. Replacing that crossing by $1+6N$ alternating
crossings retains all finite boundary labels and the whole
quotient image. A fixed four-seed outside trace, with an
integer-tangle propagation induction, gives upper four for
every $N\ge0$. Thus all four families have $\mr=b=4$.

These are infinitely many knot types, not just diagrams:
their determinants are $543+1776N$ for seeds 3145/3148 and
$537+1680N$ for 3147/3149. This follows uniformly by the
parallel-edge spanning-tree formula, using the published
determinant interpretation \cite[Lemma 3.14]{stoim}.
Separate corner-region/tree enumeration and Fox determinants
check the finite coefficients and twelve sample PDs;
pointwise permutation and outside-arc traces are independently
replayed. The all-$N$ proof is written, unformalized and
externally unreviewed. The families are not claimed mutually
disjoint or globally new, and full MRC remains open.

Joseph--Pongtanapaisan already construct classical sharp
families through finite-cycle quotient labellings with
periodic twist-box rules \cite{jp}. Section 3 of their
February 2023 preprint, especially Construction 3.5 and
Corollary 3.12, has now been read; the published 2025 title
and metadata are checked separately. Full published proof
comparison and membership of our seeds in their families
remain open. This source must be credited when discussing
periodic finite-quotient constructions.

The known-family audit also retains a specific limitation:
the published Stoimenow discussion distinguishes an
arborescent link from a link with an arborescent alternating
diagram \cite[Proposition 6.3, Remark 6.4]{stoim}. Failure
of series-parallel reduction on one checkerboard graph
alone is insufficient for knot nonmembership. No new
worldwide novelty conclusion is registered.

\section{Additivity of the strongest simple-cover bounds and a strict limit}

The working proof \path{SIMPLE_COVER_RANK_ADDITIVITY.tex}
defines $C(K)=\max_{r,\phi}(r-1+\operatorname{rank}\pi_1Y_\phi)$ and,
for a fixed field $k$, $A_k(K)=\max_{r,\phi}(r-1+b_1(Y_\phi;k))$.
The covers are connected and meridians are transpositions.
Both maxima are finite, since $r\le b(K)+1$, and
$A_k\le C\le\mr\le b$.

The written proof establishes, for all classical knots,
\[
 C(K_1\#K_2)=C(K_1)+C(K_2)-1,\qquad
 A_k(K_1\#K_2)=A_k(K_1)+A_k(K_2)-1.
\]
The same field must be used on both sides. A knot group is
normally generated by one meridian. Therefore a restricted
transposition coloring has one active sheet orbit and fixes
the complement pointwise. If the summand supports have sizes
$r_i$, share the moving pair, and have union of size $R$,
the lifted decomposing spheres give
\[
 Y\cong\#_iY_i\ \#\ \#^{\beta}(S^1\times S^2),\qquad
 \beta=\sum_i(r_i-2)-(R-2).
\]
Ordinary rank adds by the classical Grushko theorem
\cite[Theorem 1.B.5]{stallings}; fixed-field first homology
adds directly. Classifying every restriction proves the
upper maximization inequality, while taking disjoint
additional sheets realizes maximizing summand covers.

Consequently knots with $C=\mr$ satisfy meridional additivity
under connected sum, and matching bridge certificates remain
sharp. In particular, a sum of $t$ members of the four
periodic families and $s$ trefoils has $\mr=b=3t+s+1$.
Schubert bridge additivity, with primary modern proof
\cite[Theorem 1]{schultens}, makes the nonnegative defects
$b-C$ and $b-A_k$ additive as well. Earlier conditional
meridional additivity through Coxeter quotients is explicitly
credited to the introduction of \cite{bbk}.

This does not make simple-cover rank universal. The known
torus knot $T(4,7)$ has $C=3<\mr=b=4$. Every simple coloring
has degree at most five; the central torus relation excludes
onto $S_3,S_4,S_5$ images, since none has an element of order
seven. The double-cover group has an exact two-generator
presentation, cyclic homology $\mathbb Z/7$, and a noncommuting
binary matrix image. Thus $A_k=2$ at characteristic seven
and $A_k=1$ otherwise. For $J_n=\#^nT(4,7)$,
\[
 C(J_n)=2n+1,\qquad
 A_k(J_n)=1+n[\operatorname{char}k=7],\qquad
 \mr(J_n)=b(J_n)=3n+1.
\]
The last equality is a known graph-manifold case \cite{bdjw}
together with bridge additivity, not a new MRC claim. These
knots also have no $S_3$ transposition coloring, excluding
the rank-two-form part of the original binary criterion in
any dimension; higher-rank forms are not thereby excluded.

Eight literal connected-sum diagrams check the sphere graph,
folded relators and integral homology, with independent arc
tracing, reverse-word reconstruction and Smith-step replay.
The torus boundary has separate Schreier-word and matrix
checks. Four corruptions reject. Uniform normal-generation,
topological gluing and Grushko arguments remain written proofs,
unformalized and externally unreviewed. General meridional
additivity and full MRC remain open; worldwide priority of
this cover-bound formulation is not claimed.

\section{Double-transposition meridians beyond binary transvections}

For a transitive degree-$r$ permutation quotient with meridian
index $\delta=r-\#\{\text{cycles}\}$, the standard edge count gives
$\mr\ge\lceil(r-1)/\delta\rceil$: any $d$ meridians factor into
$d\delta$ transposition edges, which must connect all sheets.
Double-transposition meridians have $\delta=2$. With the retained
upper four, the only transitive degrees giving this lower four
are eight and nine. Both are now excluded on the original
sixteen PDs. The frozen searches test respectively $18\,450$
and $6\,561$ transitive four-seed orbits. Independent full
orbit marking covers all $9\,261\,000$ and $54\,010\,152$
normalized raw tails, followed by geometric arc and pointwise
relator replay. No hit occurs; enhanced cover-rank bounds and
other meridian cycle types are not excluded.

The positive control $T(4,15)$ has full $A_8$ image, with
$X=(0\ 7\ 5\ 1)(2\ 6\ 4\ 3)$,
$Y=(0\ 1\ 2\ 3\ 4)(5\ 6\ 7)$ and
$X^{-1}Y^4=(0\ 6)(2\ 5)$. The same matrices certify the known
family $T(4,60j+15)$ for every $j\ge0$, using the standard
meridian $x^{-1}y^{15j+4}$ and the four-strand upper bound.
This is not a new torus-knot MRC theorem. Finite-cycle quotient
methods already have the precedents \cite{bk,jp}.

The full image has order $20160>1344$, outside binary affine
three-space. Although $A_8\cong\mathrm{GL}_4(2)$ \cite{atlasa8},
the meridian classes differ: the double-transposition class
has size $210$, versus $105$ binary transvections $I+vf$.
An isomorphism preserves conjugacy-class size. This precisely
distinguishes the new test from the pinned external binary
rank-one search without questioning its valid certificates.
The edge-count argument above gives the uniform bound; the former
permutation-index working note is not distributed in this checkout.
These double-transposition tests do not settle the sixteen targets.
Their equalities are established by the retained graph-cover proofs.
Full MRC remains open, and worldwide novelty is unclaimed.

\section{Complete ternary affine bases and new cover quotients}

The next finite field gives a complete three-dimensional
rank-one base test. In $\mathrm{GL}_3(\mathbb F_3)$ the two
possible classes are $104$ order-three transvections and
$117$ order-two reflections. The frozen searches exhaust
$21\,632$ and $34\,084$ four-seed conjugacy orbits on the
original sixteen PDs, including all proper images.
Independent raw orbit marking covers $1\,124\,864$ and
$1\,601\,613$ tails; geometric signed traces and affine-point
coefficient replay agree with the producer.

Order three gives only constant cyclic bases. Reflections
give two full $\mathrm{GL}_3(\mathbb F_3)$ representation
orbits on each of the first eight targets, in addition to
the constant base; the last eight have only the constant
base. Root dimensions are respectively three and one, equal
to the principal-subspace dimensions. Thus every admissible
affine lift of each rank-one base is principal. A trivial
linear base gives only cyclic translations. Consequently
this ternary affine-three-space architecture cannot give a
transitive27point or homogeneous4 fixed-space certificate
on the sixteen targets. Bases over larger fields and higher
dimensions are outside this exclusion.

Both meridian classes have explicit positive controls on
the known sum of three trefoils. Four homogeneous matrices
have rank-one differences, satisfy adjacent braid relations
and have zero common fixed space. Their affine image orders
are $151632$ and $648$, with transitive27point actions;
literal nine-crossing PDs and four-seed uppers independently
give $\mr=b=4$. These are known composite controls.

The sixteen nonconstant reflection representations induce
full $\mathrm{SL}_3(\mathbb F_3)$ quotients of the double-cover
groups. Ordinary natural3 and dual3 cocycles, with no
meridian extension constraint, all have $\dim Z^1=4$ and
$\dim H^1=1$, below the sharp threshold $\dim Z^1\ge7$.
This tests32modules, not all modules or fields, and does not
determine actual cover rank. Representation orbits are not
asserted to have distinct quotient kernels.

The thirteen-point projective action gives a distinct next
augmentation gate: a homology gain at least13 at $p\ne13$
forces cover rank three. On the known control
$G_0=C_3*C_3*C_3$, the explicit quotient has point-stabilizer
preimage $F_6*(*_{12}C_3)$, with gain $18-3=15$ over
$\mathbb F_3$. Finite cycle inputs pass; subgroup structure
uses the classical graph-of-groups theorem \cite{serre}.
That checkpoint is now completed by the thirteen-point audit
below. The former ternary-affine working note and its exhaustive search
inputs are not distributed in this checkout. This test does not determine
the targets' cover ranks; the retained graph-cover proofs establish their
equalities. Full MRC remains open.

\section{Exact thirteen-point homology and the characteristic-thirteen boundary}

The rank implication of the Betti-gain argument extends to
every characteristic. For any index-$r\ge2$ subgroup of a
finitely generated group and any field,
\[
 \operatorname{rank}G\ge1+\left\lceil
       \frac{b_1(H)-b_1(G)}{r-1}\right\rceil.
\]
Indeed $\dim Z^1(G,k[G/H])=b_1(H)+r-1$ by Shapiro, and
the cocycle sum map to $Z^1(G,k)$ has rank at most $b_1(G)$.
Its kernel is the augmentation cocycle space, whose
dimension is at most $(r-1)\operatorname{rank}G$. No splitting is needed.
The earlier index-five rank criterion therefore works at
five too, although its split cohomology identity need not.
A positive degree13/characteristic13 control on
$C_{13}*C_{13}*C_{13}$ has subgroup $*_{26}C_{13}$,
gain23 and augmentation $Z^1$ dimension36, whereas the
split identity would give35. The same construction works
for every odd prime, by explicit Schreier rewriting.

The32 point/dual actions induced by the new ternary maps give
sixteen inequivalent covers, two per source knot. Their
integral homology is
$\mathbb Z^2\oplus C_2\oplus C_{882878}$ for
$13a_{129},13a_{136},13a_{194},13a_{197}$ and
$\mathbb Z^2\oplus C_2\oplus C_{3128602}$ for
$13a_{135},13a_{137},13a_{196},13a_{198}$.
Base homology is respectively $C_{391}$ and $C_{401}$.
The maximum Betti gain at $p\ne13$ is four, attained at two,
below the required thirteen. Thus the dimension12
augmentation module has ordinary cocycle dimension at most
sixteen away from thirteen. Direct characteristic13
coefficient matrices have kernel dimension fourteen, below
the required twenty-five. This completes the specified
induced-cover augmentation test over every field.

Explicit sheet intertwiners identify the second retained
representation's point action with the first one's dual
action. Both projective actions are faithful, so there is
one quotient kernel per knot among these maps. Point and
dual actions are inequivalent: the enumerated point
stabilizer432 fixes no dual point. No cross-knot pooling or
exhaustion of every index13 subgroup or quotient is inferred.

The earlier audit reported separate Smith/augmentation replay and known
positive controls. The former projective-thirteen working note and full
replay inputs are not distributed in this checkout, so these finite
observations are not a current runnable certificate.

These homology tests do not decide actual cover rank. The retained
structural proofs below establish all sixteen equalities; full MRC
remains open.

\section{Two common double covers from exact mutation certificates}

The eight sources with new ternary quotients now form two certified
mutation blocks:
\[
 \{13a_{129},13a_{136},13a_{194},13a_{197}\},\qquad
 \{13a_{135},13a_{137},13a_{196},13a_{198}\}.
\]
Their pinned thirteen-crossing reduced alternating diagrams have
seven-vertex selected Tait graphs. Six root-to-member comparisons have
explicit graph2isomorphism certificates, requiring at most one Whitney
switch before an abstract graph isomorphism. Independent corner-region
reconstruction, spanning-tree enumeration and the full binary cycle
spaces pass, with four mathematical corruptions rejected. The selected
checkerboard crossing conventions agree within each block.

Greene's Corollary4.5 supplies the graph-to-mutation implication;
only its elementary Section4 argument and explicit classical graph
inputs are used. Morton's Theorem3 gives the short torus-gluing proof
that mutation preserves double branched covers, crediting Viro1976.
Neither the full Greene lattice/Floer proof nor Viro's original full proof
has been audited. The source identities and stated scope remain in
\texttt{sources.json}; the separate mutation-source audit is not
distributed in this checkout.

Thus these eight ordinary cover-rank questions reduce to two groups,
with respective abelianizations $C_{391}$ and $C_{401}$. The mutation
argument alone bounds each rank between two and three. The four-piece
proof below establishes rank three for both, settling all four knots in
each block using their individual four-bridge uppers and the double-cover
inequality. No bridge or meridional-rank mutation invariance is assumed.
The mutation proof is
\texttt{MUTATION\_COMMON\_DOUBLE\_COVERS.tex}; the other eight original
targets use the direct three-piece proofs below.

\section{Compact core presentations and a common four-port assembly}

The two common cover groups now have exact short presentations on three
generators with relator lengths $(9,13,15)$ and $(12,15,11)$. Every
Tietze elimination, free automorphism with inverse, and relator-subword
replacement is retained and independently replayed. Their integral
abelianizations remain $C_{391},C_{401}$, and transported projective
permutations still give nonabelian quotients of order5616. Compact
presentation alone does not determine rank; the four-piece argument below
gives rank three for both groups.

Cutting each source into the first seven and last six crossings gives
four literal ports. Both inner formal normalized core groups reduce
to the same two-generator one-relator group $P$ of relator length11;
both outer ones to the same group $Q$ of relator length7. The two closed
groups are reconstructed exactly as quotients of
$P*Q*\langle z\rangle$ by four displayed port identifications. A free
base change and seventeen-to-thirteen arc identification prove this
algebraic assembly. No injective torus amalgam, geometric piece
classification or JSJ theorem is inferred. The canonical proof and
literal presentations are in
\texttt{COMPACT\_DOUBLE\_COVER\_PRESENTATIONS.tex}.

Outer-fragment propagation stalls on five inner arcs, so its two-generator
presentation gives no rank-two closed-group certificate. Abstract
series/parallel reductions give tetrahedral graph cores and branch
tree/forest counts $(19,14)$ and $(14,19)$; these conductances have not
been certified as rational-tangle filling slopes. Candidate filled
triangulations have degenerate numerical shapes and neither target is
recognized in a single frozen Regina pass. These are explicit source
and proof obligations, not graph-manifold or hyperbolicity conclusions.
Nine altered compact/assembly inputs are rejected. The new work adds
no source-catalogue entry or target equality; the catalogue remains42.
Brunner/Ito and Boileau--Weidmann papers are leads for a subsequent
peripheral/gluing audit, not imported rank classifications here.

\section{Eight sharp cases from four canonical Seifert pieces}

The two shared cover roots now have group rank exactly three.
The formal pieces from the preceding section admit exact peripheral
models over a disk with cone orders2,3 and an annulus with one cone
of order2. Their four port relations reduce to an injective torus
gluing. Thus each source group is the fundamental group of a constructed
closed orientable irreducible graph manifold with canonical chain
\[
 D(2,3)\;---\;A(2)\;---\;A(2)\;---\;D(2,3).
\]
All three maps are integral unimodular, and adjacent fiber intersections
are one. Negative orbifold Euler characteristics exclude ambiguous
fiber exceptions; the classical JSJ criterion gives four canonical pieces.
No homeomorphism from a numerical source triangulation is used.

Boileau--Weidmann's necessary rank-two classification, imported from
their Theorem2, permits at most three JSJ pieces. Therefore both source
groups have rank at least three, and their exact three-generator uppers
give equality. The double-cover bound and each member's independently
replayed four-seed upper then prove
\[
 \mr(K)=b(K)=4\quad
 (K\in\{13a_{129},13a_{135},13a_{136},13a_{137},
        13a_{194},13a_{196},13a_{197},13a_{198}\}).
\]
The canonical proof is \texttt{FOUR\_PIECE\_DOUBLE\_COVERS\_SHARP\_FOUR.tex}.
It includes explicit basis changes, peripheral injections, all matrices,
and the passage from group rank to knot equality. The independent
checker uses cyclic-amalgam normal forms and rejects five further
corruptions. Classical topology and the rank theorem are written/imported
arguments, not consequences of a green finite receipt. Formal verification
and external specialist review remain absent.

The three new scoped reads are Boileau--Weidmann2005, Hatcher's
Seifert/JSJ notes, and Millichap--Worden2016, bringing the catalogue to45.
The latter's Corollary5.4 eliminates the original exceptional
irregular-hyperbolic-cover case in the rank-two list. Its short deduction
is read as current context; the full hidden-symmetry proof is not audited
and the eight-case piece bound already follows from the older necessary
list. The original
Jaco--Shalen/Johannson theory is credited through the latter exposition.
The broad rank classification is classical. These eight equalities and
the other eight original targets are established by the retained case
proofs, and their graph-cover family membership follows from the final
rigidity argument below. The current retained293 panel has twenty named
equalities and273 names without a certificate here. Worldwide priority of
the explicit calculations, full classical MRC and exhaustive literature
remain open.

\clearpage
\section{Three further cases from the actual central filling slopes}

The next small frozen selection gives three more working equalities:
\[
 \mr(13a_{233})=b(13a_{233})=4,\quad
 \mr(13a_{239})=b(13a_{239})=4,\quad
 \mr(13a_{253})=b(13a_{253})=4.
\]
All three double-cover groups have rank three. The exact source/port
identities construct closed orientable irreducible graph manifolds with
canonical chain $D(2,3)-A(m)-D(2,3)$, with $m=2,3,2$ respectively.
This count alone does not exclude rank two: BW05 Theorem2 case6 allows
three pieces. Its specific necessary condition requires the middle piece
to be a two-bridge link exterior with the adjoining end fibres glued to
link meridians.

In the middle outer boundary basis $(c,f)$, where $f=s^m$ is the fibre,
the first end fibre has slope $(r,q)=(3,-2),(-2,-1),(-3,2)$.
The associated filling group has a quotient
\[
 C_m*C_{|r|}\twoheadrightarrow S_3.
\]
It therefore cannot be the infinite cyclic group of a solid torus.
But a two-component two-bridge link has unknotted components, and filling
either link meridian leaves a solid torus. This contradicts the only
three-piece rank-two type. Boundary injection, unique negative-base
fibres and nonparallel slopes certify that the constructed decomposition
is actually canonical. A source group isomorphism suffices; no numerical
cover recognition or rank-equals-genus inference is needed.

The complete written proof is
\texttt{THREE\_PIECE\_DOUBLE\_COVERS\_SHARP\_FOUR.tex}.
Independent replay checks five full core reductions, fourteen fragment
reductions, every selected port in complete amalgam normal form, the
unimodular gluing matrices, three finite filling quotients and individual
upper-four traces. Five corrupted inputs reject. The case6 source read
is separately sealed; the catalogue remains45.
Five proposed mutation transfers all failed and have unequal determinant
pairs $403/389$, $361/359$, $361/337$, $361/311$, $337/311$.
No equality is transferred to those targets. Similar PD prefixes are
not a common-cover certificate.

This three-piece argument treats three original targets. The five other
direct cases
$13a_{234},13a_{242},13a_{243},13a_{244},13a_{265}$ are covered in the next
section. All sixteen belong to the known graph-cover family by the final
rigidity argument. The working proofs and classical imports remain
unformalized and externally unreviewed; named global priority,
exhaustive source audits and full MRC remain open.

\section{Completion of the original sixteen through five direct models}

The five remaining original targets now have working proofs of
\[
 \mr(K)=b(K)=4,\qquad
 K\in\{13a_{234},13a_{242},13a_{243},13a_{244},13a_{265}\},
\]
with double-cover group rank three. These use each literal source PD
directly; the failed transfers above are retained and unused.
The exact three-piece model is again $D(2,3)-A(m)-D(2,3)$.
The middle piece has $C=\langle k,s\mid[k,s^m]=1\rangle$, fibre
$f=s^m$ and outer section $c=sk$ or the conjugate marking $ks$.
Both boundary marking order and the $h_Q^{-1}$ attachment at $265$
are checked explicitly. The actual first end-fibre filling slopes in
$(c,f)$ are
\[
 (m;r,q)=(2;3,-2)\text{ for }234,\quad
 (3;-2,1)\text{ for }242,\quad
 (2;-3,2)\text{ for }243,244,265.
\]
All fillings have a nonabelian $C_m*C_{|r|}$ quotient onto $S_3$.
This contradicts the central two-bridge meridian requirement in BW05
Theorem2 case6, using the same explicit realization, canonical JSJ and
unknotted-component filling argument as the preceding section.

The complete written proof is
\texttt{FINAL\_FIVE\_DOUBLE\_COVERS\_SHARP\_FOUR.tex}.
Independent replay verifies eight full core and thirty-two fragment
reductions, all balanced four-port cuts and source joins, exact root
eliminations, complete amalgam normal forms, all ports and conjugations,
primitive gluing matrices and five nonabelian finite filling images.
The three preceding sharp roots reproduce their sealed models, each
of the five new knots has its own checked upper-four trace, and six
corrupted inputs reject. The genuine meridian-filling control stays
cyclic. Source catalogue45 is unchanged; no rank conclusion is inferred
from piece count or a failed quotient test alone.

This completes the original sixteen-case panel here, giving twenty
named equalities in the retained293 panel and273 names without a
certificate in this programme. This is not a worldwide open/closed
census. All proofs remain unformalized and externally unreviewed. The
known-family audit is supplied by the rigidity argument below; the
sixteen are members of the existing BJW18 graph-cover family. Named global
priority, full MRC and exhaustive literature remain open. A general
extension beyond that family needs a separately justified new mechanism,
with the completed original panel retained as positive controls.

\section{Actual-cover rigidity resolves the sixteen-case known-family audit}

The original sixteen belong to an existing MRC family. Their exact group
models were already sufficient for the cover-rank proofs. A further
classical step now gives actual-cover topology: AFW15
Theorem2.1.3 classifies closed orientable manifolds with isomorphic
fundamental groups up to orientations of prime factors and lens-space
ambiguity \cite{afw}. The constructed model is closed, irreducible and
prime, with an infinite injected trefoil group. It has no nontrivial lens
summand and only one prime factor. Hence the source double branched cover
is homeomorphic to this graph manifold, possibly reversing orientation.
This handles source-cover primeness without assuming it from a table
label. No branch-involution transport or numerical recognition is needed.

The actual covers therefore satisfy the prime graph-manifold hypothesis
of the published BJW18 Theorem1.2 and Corollary1.3 \cite{bjw}.
With the certified exact bridge four, the MRC equality falls in that
2018 theorem. In fact the published rank-one/two equivalences and rank-three
graph-cover theorem already imply MRC throughout this cover class whenever
$b\le4$: decide the meridional rank among1,2,3,4 and apply the corresponding
bridge consequence. The independent rank-three calculation remains useful
because it determines the exact common value four.

The complete source-scope proof is
\texttt{ORIGINAL\_SIXTEEN\_KNOWN\_GRAPH\_COVER\_FAMILY.tex}.
Sixteen literal PD/upper bindings, the three retained lower/model receipts,
and the exact mutation aliases for the first eight are sealed separately.
The finite binder does not prove rigidity or reprove the published MRC
theorem. The AFW final author book Theorems2.1.2--2.1.3 and their earlier
author-v3 counterparts were read at the precise lens/orientation scope;
the source pages were visually checked. The current catalogue contains46
sources.

This establishes the known-family scope for all sixteen while preserving
every equality and exact finite certificate. They are members of a known
MRC class. Named global bridge/rank computational priority
is still unaudited. The other four sharp simple-three-cover seeds require
their own family audit. A prospective new target outside BJW18 needs
actual double-cover topology outside that class; a hyperbolic knot exterior
or non-arborescence alone does not suffice. Full MRC, exhaustive literature,
formal and external specialist review remain open.

{\small

\section{Arbitrary-cycle covers: twenty-four further sharp equalities}

The current proof \path{GENERAL_BRANCHED_COVER_RANK_BOUNDS.tex} retains
the fundamental group of a branched cover with arbitrary meridian
cycle type. For a transitive degree-$n$ permutation representation,
let $c$ be the number of cycles of a meridian, including fixed points,
and let $Y$ be the branched completion. Eliminating one lifted edge
per filled meridian cycle gives
\[
 \mr(K)\ge
 \left\lceil\frac{n-1+\operatorname{rank}\pi_1(Y)}{n-c}\right\rceil.
\]
Its homology version is valid over every field.
The inequality strictly improves the maximal simple-cover rank
invariant on the known family $T(4,420j+7)$: $C=3<D=4$ for all $j\ge0$.
This is a comparison of methods on known torus knots, without a
priority claim for the general bound.

The retained twenty-nine binary $GL_3$ witnesses now have a different
use from their deficient rank-one coefficient tests. Their natural
and dual seven-point actions have meridian index two. Exact group
presentations of all fifty-eight branched covers reduce by explicit
one-occurrence Tietze eliminations: forty-eight are $C_3$, and ten
are trivial. The former give $\mr=b=4$ on the twenty-four names in
the general-cover theorem, with separate four-seed bridge upper
certificates. A full directed-edge computation independently gives
$b_1(Y;\mathbb F_3)=1$ there. The five negative names have trivial
cover groups for these fixed actions, so nonlinear quotients of these
particular covers cannot improve their bound.

Every literal PD has a current exact crossing-preserving map to the
independently packaged named Hoste--Thistlethwaite census. All
twenty-four positive names remain in Bhattacharjee's current
293-name open list, and exact read-only row extraction from the
pinned BKM workbook finds no recorded maximal finite $S,D,H$
Coxeter hit. These comparisons exclude those recorded certificates,
not every existing geometric or infinite Coxeter family.
The small current input and standard-library checkers are retained;
no broad historical quotient search is needed to replay these proofs.

The separate \path{SEED_COVER_TOPOLOGY.tex} establishes that all
twenty-four new examples and the four earlier simple-three-cover
seeds lie outside BJW18's graph-manifold double-cover hypothesis.
Each actual double cover is irreducible and has an incompressible
torus cutting off a certified cusped hyperbolic census manifold.
The acceptance uses exact triangulation signatures and trusted
normal-surface algorithms, with an explicit imported certified-census
theorem. Numerical volume is not used. This settles the precise
BJW18 comparison on these finite seeds; global priority, fishnet and
other-family comparisons, and every positive periodic parameter
remain separate obligations.

Finally, \path{CONJUGACY_RESTRICTED_QUOTIENT_CAPACITY.tex} gives a
meridian-sensitive finite quotient gate: eight distinct
$GL_3(\mathbb F_2)$ knot-group quotient kernels with involutory
meridians force meridional rank at least four. Exact enumeration gives
only seven generating-involution-triple automorphism orbits, versus
57 unrestricted generating-pair orbits. The corresponding
double-cover statement requires an invariant kernel and inner deck
action. It supplies no additional named equality by itself.

\begin{thebibliography}{99}
\setlength{\itemsep}{2pt}
\small
\setlength{\baselineskip}{11pt}
\setlength{\itemsep}{0pt}
\bibitem{mw16} C. Millichap and W. Worden, Hidden symmetries and
commensurability of 2-bridge link complements, Pacific J. Math.285(2)
(2016), Corollaries5.3--5.4.
\url{https://doi.org/10.2140/pjm.2016.285.453}.
\bibitem{bw05} M. Boileau and R. Weidmann, The structure of 3-manifolds
with two-generated fundamental group, Topology44 (2005), 283--320.
Read author manuscript, Theorem2 and relevant Section3.1 argument.
\url{https://doi.org/10.1016/j.top.2004.10.008}.
\bibitem{hatcher} A. Hatcher, Notes on Basic 3-Manifold Topology,
online notes with revisions1999/2000, Sections1.2 and2.1.
\url{https://pi.math.cornell.edu/~hatcher/3M/3M.pdf}.
\bibitem{greene} J. E. Greene, Lattices, graphs, and Conway mutation,
Invent. Math.192 (2013), 717--750. Author preprint arXiv:1103.0487v1,
Section4.1--4.2, especially Corollary4.5.
\url{https://doi.org/10.1007/s00222-012-0421-4}.
\bibitem{morton} H. R. Morton, Mutant knots, author manuscript,
16March2014, Theorem3, p.6, crediting O. Ya. Viro (1976).
\url{https://www.liverpool.ac.uk/~su14/papers/mutantchapter.pdf}.
\bibitem{stallings} J. Stallings, Group Theory and Three-Dimensional
Manifolds, Yale University Press (1971), Theorem 1.B.5, p.~6.
\url{https://webhomes.maths.ed.ac.uk/~v1ranick/papers/stallings1.pdf}.
\bibitem{schultens} J. Schultens, Additivity of bridge numbers of
knots, Math. Proc. Cambridge Philos. Soc. 135 (2003), 539--544,
Theorem 1, crediting Schubert (1954).
\url{https://doi.org/10.1017/S0305004103006832}.
\bibitem{stoim} A. Stoimenow, Graphs, determinants of knots and
hyperbolic volume, Pacific J. Math. 232 (2007), 423--451.
\url{https://doi.org/10.2140/pjm.2007.232.423}.
\bibitem{jp} J. Joseph and P. Pongtanapaisan, Meridional rank and
bridge number of knotted 2-spheres, Canad. J. Math. 77 (2025),
282--299 (online 2023).
\url{https://doi.org/10.4153/S0008414X23000883}.
Classical construction read in Section 3 of
\url{https://arxiv.org/html/2111.09233v2}; full published proof
not independently compared.
\bibitem{serre} Serre, Trees, translated by Stillwell, Springer
(1980), Chapter I, Trees and Amalgams.
\url{https://doi.org/10.1007/978-3-642-61856-7}.
\bibitem{atlas} ATLAS of Finite Group Representations, Alternating
group $A_5$: complete ordinary and modular representation lists.
\url{https://brauer.maths.qmul.ac.uk/Atlas/alt/A5/}.
\bibitem{atlasgl} ATLAS of Finite Group Representations,
$L_2(7)=L_3(2)$: automorphism and characteristic-two simple lists.
\url{https://brauer.maths.qmul.ac.uk/Atlas/lin/L32/}.
\bibitem{atlasa8} ATLAS of Finite Group Representations,
$A_8=L_4(2)$: group order and representation list.
\url{https://brauer.maths.qmul.ac.uk/Atlas/v3/alt/A8/}.
\bibitem{bc} Burness--Crestani, On the generating graph of direct
powers of a simple group, J. Algebraic Combinatorics 38 (2013),
329--350. \url{https://doi.org/10.1007/s10801-012-0405-x}.
Author manuscript p.~2:
\url{https://seis.bristol.ac.uk/~tb13602/docs/BC_1.pdf}.
\bibitem{ps} Poudel--Saveliev, Link homology and equivariant gauge
theory, Algebraic \& Geometric Topology 17 (2017), 2635--2685,
Section 7.3. \url{https://doi.org/10.2140/agt.2017.17.2635}.
\bibitem{k3} Baykur--Kirby--Ruberman, eds., K3: A New Problem List in
Low-Dimensional Topology (2026), Problem 1.18.
\url{https://bpb-us-e2.wpmucdn.com/websites.umass.edu/dist/b/22144/files/2026/04/K3-problem-list-watermarked.pdf}.
\bibitem{rz} Rost--Zieschang (1987).
\url{https://doi.org/10.1112/jlms/s2-35.3.551}.
\bibitem{bjw} Boileau--Jang--Weidmann (2018).
Meridional rank and bridge number for a class of links,
Pacific J. Math. 292(1), Theorem1.2 and Corollary1.3.
\url{https://doi.org/10.2140/pjm.2018.292.61}.
\bibitem{afw} Aschenbrenner--Friedl--Wilton, 3-Manifold Groups,
EMS (2015), \url{https://doi.org/10.4171/154}.
Read final author manuscript Theorems2.1.2--2.1.3,
printedpp30--31; earlier author-v3 Theorems2.2--2.3 compared.
\bibitem{bdjw} Boileau--Dutra--Jang--Weidmann (2017).
\url{https://arxiv.org/abs/1608.01570}.
\bibitem{cornwell} Cornwell (2014).
\url{https://arxiv.org/abs/1303.4943}.
\bibitem{cornwell17} Cornwell, KCH representations, augmentations,
and A-polynomials, J. Symplectic Geometry 15 (2017), 983--1017,
Theorem 1.10 and its proof on arXiv p.~18.
\url{https://arxiv.org/abs/1310.7526v3}.
\bibitem{ch} Cornwell--Hemminger (2016).
\url{https://arxiv.org/abs/1408.4110}.
\bibitem{wirt} Blair--Kjuchukova--Velazquez--Villanueva (2020).
\url{https://archive.intlpress.com/site/pub/files/_fulltext/journals/cag/2020/0028/0002/CAG-2020-0028-0002-a002.pdf}.
\bibitem{bbk} Baader--Blair--Kjuchukova (2021).
\url{https://doi.org/10.1007/s00208-020-02124-z}.
\bibitem{twigs} Baader--Blair--Kjuchukova--Misev (2023).
\url{https://arxiv.org/abs/2008.00763}.
\bibitem{moriah} Moriah, Heegaard splittings of knot exteriors,
Geometry \& Topology Monographs 12 (2007), 191--232,
Theorems 3.6 and 4.20 (reporting Lustig--Moriah 1993).
\url{https://doi.org/10.2140/gtm.2007.12.191}.
\bibitem{core} Silver--Traldi--Williams, Core groups,
J. Knot Theory Ramifications 33 (2024), 2450013, introduction and
Section 7. \url{https://arxiv.org/abs/2311.02048v2}.
\bibitem{bk} Baader--Kjuchukova (2020).
\url{https://arxiv.org/abs/1711.08144}.
\bibitem{bkm} Blair--Kjuchukova--Morrison (2024).
\url{https://arxiv.org/abs/2208.09032}.
\bibitem{fishnet} Blair--Kjuchukova--Pfaff (2026).
\url{https://doi.org/10.1017/S030500412610214X}.
\bibitem{plain} Blair--Kjuchukova--Pfaff (2026).
\url{https://doi.org/10.2140/pjm.2026.342.217}.
\bibitem{whitehead} Gao--Gu--Li--Wang (2026), preprint.
\url{https://arxiv.org/abs/2611.15673v1}.
\bibitem{deposit} Bhattacharjee (2026), unrefereed deposit.
\url{https://zenodo.org/records/23561157}.
\bibitem{cesy} Clark--Elhamdadi--Saito--Yeatman, Quandle Colorings of
Knots and Applications (2014), Sections 6 and 7.
\url{https://arxiv.org/abs/1312.3307}.
\bibitem{musick} Musick, Minimal bridge projections for 11-crossing
prime knots (2012), Theorem 1. \url{https://arxiv.org/abs/1208.4233v3}.
\bibitem{kl} Kauffman--Lambropoulou, On the classification of rational
tangles (2004), Theorem 1 and Lemma 2.
\url{https://arxiv.org/abs/math/0311499v2}.
\bibitem{preorder} Boileau--Kitano--Nozaki (2025), preprint.
\url{https://arxiv.org/abs/2510.12372}.
\end{thebibliography}
}
\end{document}
